\subsection{Интегрирование иррациональных выражений вида $R\!\left(x,\,\sqrt{\dfrac{ax+b}{cx+d}}\right)$}

\subsubsection*{1. Общая схема}

Рассматривается интеграл вида
\[
\int R\!\left(x,\,\sqrt{\frac{ax+b}{cx+d}}\right)dx,
\]
где $R(u,v)$ — рациональная функция двух переменных, $ad - bc \neq 0$.

\textbf{Рационализирующая замена:}
\[
t = \sqrt{\frac{ax+b}{cx+d}},
\]
то есть $t^2 = \dfrac{ax+b}{cx+d}$.

Из этого соотношения выражается $x$:
\[
t^2(cx+d) = ax+b \implies x(ct^2 - a) = b - dt^2 \implies
x = \frac{b - dt^2}{ct^2 - a}.
\]

Дифференцируя, находим $dx$ и подставляем. После замены интеграл становится \textbf{рациональным} по $t$ и вычисляется по стандартной схеме разложения на простейшие дроби.

\subsubsection*{2. Частный случай: $\sqrt[n]{ax+b}$}

Если иррациональность имеет вид $\sqrt[n]{\,ax+b\,}$, то замена $t = \sqrt[n]{ax+b}$ (т.\,е. $x = \dfrac{t^n - b}{a}$) рационализирует интеграл.

Если одновременно присутствуют $\sqrt[m]{\,ax+b\,}$ и $\sqrt[k]{\,ax+b\,}$, выбирают $n = \mathrm{lcm}(m,k)$ и полагают $t = \sqrt[n]{ax+b}$.

\subsubsection*{3. Пример: $\displaystyle\int \frac{dx}{\sqrt[3]{(x-1)(x+1)^2}}$}

Запишем подкоренное выражение:
\[
(x-1)(x+1)^2 = (x+1)^3 \cdot \frac{x-1}{x+1}.
\]
Поэтому:
\[
\sqrt[3]{(x-1)(x+1)^2}
= (x+1)\cdot\sqrt[3]{\frac{x-1}{x+1}}.
\]
Интеграл принимает вид:
\[
\int \frac{dx}{(x+1)\,\sqrt[3]{\dfrac{x-1}{x+1}}}.
\]

\textbf{Замена:}
\[
t = \sqrt[3]{\frac{x-1}{x+1}} = \left(\frac{x-1}{x+1}\right)^{1/3},
\quad t^3 = \frac{x-1}{x+1}.
\]

Выражаем $x$: $t^3(x+1) = x-1 \implies x(t^3-1) = -1-t^3 \implies$
\[
x = \frac{-(1+t^3)}{t^3-1} = \frac{1+t^3}{1-t^3}.
\]

Находим $x+1$:
\[
x+1 = \frac{1+t^3}{1-t^3}+1 = \frac{2}{1-t^3}.
\]

Дифференцируем $x$:
\[
dx = \frac{3t^2(1-t^3)-\,(1+t^3)\cdot(-3t^2)}{(1-t^3)^2}\,dt
= \frac{3t^2(1-t^3)+3t^2(1+t^3)}{(1-t^3)^2}\,dt
= \frac{6t^2}{(1-t^3)^2}\,dt.
\]

Подставляем в интеграл:
\[
\int \frac{1}{\dfrac{2}{1-t^3}\cdot t}\cdot\frac{6t^2}{(1-t^3)^2}\,dt
= \int \frac{(1-t^3)}{2t}\cdot\frac{6t^2}{(1-t^3)^2}\,dt
= \int \frac{3t}{1-t^3}\,dt.
\]

Разложим $1-t^3 = (1-t)(1+t+t^2)$:
\[
\frac{3t}{1-t^3} = \frac{3t}{(1-t)(1+t+t^2)} = \frac{A}{1-t} + \frac{Bt+C}{1+t+t^2}.
\]

Подбор коэффициентов: $t=1 \implies 3 = 3A \implies A=1$;
приравниваем коэффициенты при $t^2$: $0 = -A + (-B) \implies B = -1$;
свободные члены: $0 = A+C \implies C = -1$.

Таким образом:
\[
\frac{3t}{1-t^3} = \frac{1}{1-t} - \frac{t+1}{1+t+t^2}.
\]

Интегрируем:
\[
\int\frac{dt}{1-t} = -\ln|1-t|,
\]
\[
\int\frac{t+1}{1+t+t^2}\,dt
= \frac{1}{2}\int\frac{2t+1}{1+t+t^2}\,dt + \frac{1}{2}\int\frac{dt}{(t+\frac{1}{2})^2+\frac{3}{4}}
= \frac{1}{2}\ln(1+t+t^2) + \frac{1}{\sqrt{3}}\arctan\frac{2t+1}{\sqrt{3}}.
\]

\textbf{Ответ} (в терминах $t = \sqrt[3]{\dfrac{x-1}{x+1}}$):
\[
\boxed{
\int\frac{dx}{\sqrt[3]{(x-1)(x+1)^2}}
= -\ln|1-t| - \frac{1}{2}\ln(1+t+t^2) - \frac{1}{\sqrt{3}}\arctan\frac{2t+1}{\sqrt{3}} + C.
}
\]